This chapter presents the proposed solution: a modular parallel framework for
the Aho-Corasick multi-pattern matching algorithm, targeting shared-memory
multi-core CPUs. The design is guided by a single structural decision from which
every other property follows. The execution is split into separate phases: the
master thread builds the automaton sequentially, after which its
state-transition tables are declared strictly read-only. Because the automaton
is never mutated during scanning, multiple worker threads can traverse it
concurrently without any mutual-exclusion primitive on the hot path. This
lock-free property is what allows the parallel search to approach the limits
imposed by the hardware rather than by software contention.

The framework is organized so that distinct optimization ideas can be plugged in
and measured independently. Two complementary forms of parallelism are
considered: partitioning the \textit{input text} among threads, and partitioning
the \textit{pattern dictionary}. In addition, an algorithmic optimization of the
match-emission step is proposed that benefits both the sequential and the
parallel cases. The remainder of this chapter describes each component.

\subsecao{Data Partitioning and Boundary Overlap}

To exploit thread-level parallelism on a single large input, the framework
partitions the text buffer into contiguous segments, one per worker thread.
Each worker traverses the shared automaton over its own segment and records
matches in a private list, so that no synchronization is required while
scanning.

Partitioning a stream of bytes naively would, however, miss any pattern that
straddles the boundary between two adjacent segments, because neither worker
would observe the full pattern. To prevent these boundary false negatives, the
design uses a controlled overlap. Let $L_{\max}$ be the length of the longest
pattern in the dictionary. A worker responsible for the segment spanning offsets
$[S, E)$ does not begin scanning at $S$ but at $S - L_{\max} + 1$, except for the
first worker, which starts at offset $0$. This warm-up margin guarantees that, by
the time the worker reaches the start of its owned region at $S$, the automaton is
in the same state it would have reached in a single sequential pass from offset
$0$, so every boundary-crossing pattern is observed by exactly one worker. Any
match whose end position falls before $S$ is suppressed, because it is owned and
reported by the previous worker; conversely, attributing each match to the segment
that contains its end position keeps ownership disjoint, which preserves an exact
one-to-one correspondence with the sequential baseline and duplicates no
detection. An overlap of $L_{\max}-1$ is the minimum that achieves
this property; any smaller value loses matches, and any larger value only adds
redundant work.

\subsecao{Load Balancing on Heterogeneous Cores}

Uniform segments are optimal only when all cores are equally fast. On a hybrid
processor, where Performance and Efficiency cores run at different clock
frequencies, equally sized segments cause the threads on slower cores to finish
late, leaving the faster threads idle at the final synchronization barrier. The
proposed framework therefore includes a topology-aware partitioning strategy
that detects the heterogeneous layout and weights each segment in proportion to
the relative speed of the core that will process it. The aim is to equalize the
finishing times of all workers and thereby reduce the idle time wasted at the
barrier. As an alternative that requires no knowledge of the topology, the
framework also supports dynamic dispatch, in which the text is divided into many
small tasks that idle workers claim through a single atomic counter, so that
faster cores naturally process more tasks.

\subsecao{Flattening the Match-Emission Path}

In the canonical Aho-Corasick algorithm, each accepting state references the
patterns that end there through linked structures, and reporting all matches at
a given state requires walking these pointer chains. Under a large dictionary,
those pointers are scattered across an automaton far larger than the cache, so
each emission incurs cache misses and serialized pointer dereferences that
compete with the much hotter state-transition traffic.

To remove this cost, the design proposes a flattened output layout. During
construction, the patterns associated with every state are collected into a
single contiguous array, and each state stores only a starting offset and a
count into that array. At scan time, emitting the matches of a state becomes a
sequential read of a short, contiguous range, which the hardware prefetcher
handles efficiently and which removes pointer dereferences from the inner loop.
This optimization is purely a memory-layout change: it preserves the exact set
of reported matches and is inherited by every scanning variant, whether
sequential or parallel.

\subsecao{Partitioning the Dictionary}

The strategies above all partition the text while keeping the dictionary whole.
A complementary axis, evaluated as a contrast, partitions the dictionary
instead. Here the pattern set is divided into $K$ groups according to the first
byte of each pattern, and a smaller sub-automaton is built for each group. Every
sub-automaton is then scanned over the input. The motivation is that each
sub-automaton has a smaller working set, which might fit closer to the cache and
mitigate the memory pressure of a single monolithic automaton. The framework
also supports a two-dimensional composition that combines dictionary sharding
with text chunking. Whether either approach actually outperforms plain text
parallelism on memory-bound workloads is an empirical question, answered in the
results chapter.
